% SPDX-License-Identifier: AGPL-3.0-or-later
% Commercial license available
% © Concepts 1996–2026 Miroslav Šotek. All rights reserved.
% © Code 2020–2026 Miroslav Šotek. All rights reserved.
% ORCID: 0009-0009-3560-0851
% Contact: www.anulum.li | protoscience@anulum.li
% SCPN-FUSION-CORE — coupled transport model intersection with TORAX.
\documentclass[11pt]{article}

\usepackage[utf8]{inputenc}
\usepackage[T1]{fontenc}
\usepackage{amsmath}
\usepackage{booktabs}
\usepackage{microtype}
\usepackage{natbib}
\usepackage[colorlinks=true,linkcolor=blue,citecolor=blue,urlcolor=blue]{hyperref}
\usepackage[margin=1in]{geometry}
\bibliographystyle{unsrtnat}

\title{A Four-State Coupled Transport Step Checked Against TORAX on a Declared
Model Intersection, with End-to-End Gradients}
\author{Miroslav \v{S}otek\\
\small ANULUM CH\&LI, Marbach SG, Switzerland, \url{https://www.anulum.li}\\
\small ORCID: \href{https://orcid.org/0009-0009-3560-0851}{0009-0009-3560-0851}\\
\small \href{mailto:protoscience@anulum.li}{protoscience@anulum.li}}
\date{Preprint --- 16 November 2026}

\begin{document}
\maketitle

\begin{abstract}
We check a coupled transport step of SCPN-Fusion-Core against TORAX 1.4.3 on a
model intersection that both codes implement: circular geometry, constant
diffusivities, Gaussian heat, particle and current sources, collisional
ion--electron exchange and prescribed edge values.  The step advances the ion
and electron temperatures, the electron density and the poloidal flux together
from one beginning-of-step state.  The case is short by design: two steps of
10\,ms on 16 cells.  At the final time the two codes differ by a relative $L_2$
of $5.4 \times 10^{-4}$ in density, $3.9 \times 10^{-2}$ in electron temperature,
$1.1 \times 10^{-2}$ in ion temperature and $6.9 \times 10^{-3}$ in poloidal
flux, inside thresholds declared before the run.  Halving the time step changes
each code by at most $1.7 \times 10^{-3}$, so the difference between the codes
is not a time-step effect of either.  Source totals close to $4.4 \times
10^{-16}$ in both codes, and repeated runs are byte-identical.  A JAX version of
the same step gives gradients that agree with central finite differences to
$4.1 \times 10^{-8}$ and reproduces the production step to $7.1 \times
10^{-15}$.  A bounded optimisation of three source scales lowers a target
misfit 174-fold, but it does not recover the scales that produced the target:
over 20\,ms the driven-current scale is nearly unobservable.  Nothing here
claims equivalence with TORAX outside the intersection, experimental
validity, or a speed advantage.
\end{abstract}

\noindent\textbf{Keywords:} tokamak transport; code comparison; TORAX;
automatic differentiation; current diffusion; model intersection

\section{Introduction}

TORAX solves the coupled equations for ion heat, electron heat and particle
transport and for current diffusion, in JAX, with just-in-time compilation and
automatic differentiation \citep{Citrin2024}.  Comparing another transport
code with it is only meaningful where both implement the same physics.  We
therefore declare a model intersection, a set of models that both codes
provide, and compare the codes on it alone.  The comparison checks the
numerics of a coupled step, its source and conservation budgets, its
determinism and its differentiability.  It does not check transport physics
outside the intersection.

\section{The model intersection}

The deck fixes every input (Table~\ref{tab:deck}).  The geometry is circular
with an elongation of one, at ITER-like size; it is not an ITER scenario.  The
transport coefficients are constant: $\chi_i = \chi_e = 1$\,m$^2$\,s$^{-1}$ and
$D = 0.3$\,m$^2$\,s$^{-1}$, with no convection.  The sources are Gaussian in
the normalised radius.  The edge values change from their initial to their final values over the run.  The case
runs from 0 to 20\,ms in two steps of 10\,ms on 16 TORAX cells, and the
profiles are compared at 17 points.  A refined run uses steps of 5\,ms.  The
thresholds were fixed in the deck before the runs.

\begin{table}[ht]
\centering
\caption{The declared model intersection.}
\label{tab:deck}
\small
\begin{tabular}{ll}
\toprule
Geometry & circular, $R = 6.2$\,m, $a = 2.0$\,m, $B = 5.3$\,T, $\kappa = 1$ \\
Grid and time & 16 cells, 17 comparison points; 0--20\,ms, $\Delta t = 10$\,ms (refined 5\,ms) \\
Plasma & $I_p = 5$\,MA, $Z_\mathrm{eff} = 1.5$ \\
Initial core and edge & $T_i$ 5 and 0.5\,keV; $T_e$ 4 and 0.4\,keV; $n_e$ $5 \times 10^{19}$ and $2 \times 10^{19}$\,m$^{-3}$ \\
Transport & $\chi_i = \chi_e = 1$\,m$^2$\,s$^{-1}$, $D = 0.3$\,m$^2$\,s$^{-1}$, resistivity multiplier 20 \\
Heat source & 10\,MW, 60\% to electrons, centre 0.3, width 0.25 \\
Particle source & $10^{20}$\,s$^{-1}$, centre 0.45, width 0.25 \\
Driven current & 100\,kA, centre 0.35, width 0.2 \\
\bottomrule
\end{tabular}
\end{table}

The SCPN step advances $T_i$, $T_e$, $n_e$ and the poloidal flux $\psi$ from
the same beginning-of-step state.  It keeps explicit accounts of the sources,
the boundary values, the exchange and the linear solves.  The ordinary
transport solver of the code remains its production path; this step is a
separate surface for code-to-code comparison.  TORAX 1.4.3 ran on the CPU in
float64 through a pinned reference runner.

\section{Agreement with TORAX}

Table~\ref{tab:parity} lists the final-state differences.  All four lie
inside their thresholds.  The largest, $3.9 \times 10^{-2}$ in the electron
temperature, is about a third of its threshold of 0.12.

\begin{table}[ht]
\centering
\caption{Relative $L_2$ difference between the codes at the final time, the
declared threshold, and the change of each code when the time step is halved.}
\label{tab:parity}
\small
\begin{tabular}{lrrrr}
\toprule
 & SCPN vs TORAX & Threshold & SCPN refined & TORAX refined \\
\midrule
$n_e$ & $5.44 \times 10^{-4}$ & 0.08 & $3.29 \times 10^{-5}$ & $2.95 \times 10^{-4}$ \\
$T_e$ & $3.88 \times 10^{-2}$ & 0.12 & $2.64 \times 10^{-4}$ & $1.09 \times 10^{-3}$ \\
$T_i$ & $1.12 \times 10^{-2}$ & 0.12 & $2.54 \times 10^{-4}$ & $1.72 \times 10^{-3}$ \\
$\psi$ & $6.94 \times 10^{-3}$ & 0.20 & $2.32 \times 10^{-5}$ & $1.35 \times 10^{-4}$ \\
\bottomrule
\end{tabular}
\end{table}

Halving the time step changes each code by at most $1.7 \times 10^{-3}$, well
below the declared 0.02.  The difference between the codes in $T_e$ is about
35 times the larger of the two refinement changes, so it does not come from
either code's time discretisation.  Its source within the intersection, for
example the spatial discretisation or the treatment of the exchange term, is
not identified here.

The budgets close to rounding.  The source totals agree with their declared
integrals to $1.5 \times 10^{-16}$ in SCPN and $4.4 \times 10^{-16}$ in TORAX,
against a threshold of $5 \times 10^{-3}$.  In SCPN the ion--electron exchange
closes to $1.6 \times 10^{-17}$ and the largest linear residual is $7.1 \times
10^{-15}$.  Five replays of the SCPN step are byte-identical, and so are the
TORAX replays.

The warm wall time of TORAX is the minimum of five replays, 0.0934\,s (the
replays range from 0.0934 to 0.110\,s); SCPN's is 1.50\,ms, and the declared
limit on their ratio is 100.  This gate only checks that the comparison is
complete at a bounded cost.  The two codes differ in language, compilation
and work per step, and the timings come from one loaded workstation, so they
support no speed comparison.

\section{End-to-end gradients}

A JAX version of the same step, with JAX-native tridiagonal solves, is
attached to the production solver and checked against it.  Over the two steps
it reproduces the production states with a largest absolute difference of
$7.1 \times 10^{-15}$.  Its gradients
with respect to the scales of the heat source, the particle source and the
driven current agree with central finite differences (step $10^{-4}$) to a
relative $4.1 \times 10^{-8}$ at most.  A first-order prediction of the change
of the objective under a small control change is within 1.0\% of the actual
change.

For the optimisation, a target is produced by the production step with the
three scales set to 1.25, 0.8 and 1.3.  Starting from one for each, a bounded
optimisation lowers the misfit to that target from $6.15 \times 10^{-6}$ to
$3.53 \times 10^{-8}$ in 12 iterations, a factor of 174, and the run is
byte-identical when repeated.  It does not recover the scales: it ends at
1.269, 0.872 and 0.999.  The gradients explain why.  At the starting point the
gradients with respect to the particle and driven-current scales are about
$1.4 \times 10^{4}$ and $1.6 \times 10^{5}$ times smaller than the one with
respect to the heat scale, so over 20\,ms those sources barely affect the
misfit.  The objective is identifiable here; the controls are not.

\section{Limitations}

\begin{itemize}
\item The intersection is narrow by design: constant coefficients, Gaussian
sources and circular geometry.  Nothing is claimed for the production
transport models or for TORAX outside it.
\item The case covers 20\,ms in two steps.  Longer runs and finer grids are not
compared.
\item The cause of the remaining electron-temperature difference is not
identified.
\item No experimental data enter.  Timings come from one loaded workstation
and support no speed comparison.
\end{itemize}

\section*{Software version and date}

This preprint describes \texttt{scpn-fusion} at commit \texttt{1445bb1405ed}
of 16 November 2026, on the main line after release 4.1.0.  The parity report
was produced at commit \texttt{849b9b2cfd1f} and the gradient report at
commit \texttt{d84af5398f4a}.  Every solver, benchmark and deck source those
reports record is byte-identical at the commit above.  TORAX 1.4.3 ran on the
CPU in float64 with Python 3.12.3.

\section*{Data and code availability}

Code is available at \url{https://github.com/anulum/scpn-fusion-core}.  The two
reports and the TORAX reference run cited here, with their source revisions,
Git blob identities and SHA-256 digests, are in \texttt{evidence/} of this
submission directory.  The generator in that directory checks every quoted
value against them.  This is a venue-neutral preprint and has not been
submitted.  No DOI has been reserved and the manuscript-content licence has not
yet been selected by the author.

\section*{Use of AI tools}

Large language model coding agents were used substantively for generating and
debugging the analysis code and for drafting and revising the manuscript text.
The agents were commercially available models, which are not named because no
vendor is involved in or sponsors this work.  They were coordinated through
SYNAPSE CHANNEL (version 0.99.52 at the time of writing;
\url{https://github.com/anulum/synapse-channel}), the open-source
agent-coordination tool of the ANULUM Agentic Coordination, Assurance and
Continuity Systems group.  The author specified the comparison, verified the
code through its tests and every reported number against the committed
reports, and takes full responsibility for the content.

\bibliography{references}
\end{document}
